\subsection{Síntese dos objetivos e contribuições}
\label{sec:cap5_sintese_obj_contrib}

O objetivo geral desta pesquisa foi formular e simular computacionalmente a dinâmica de uma espaçonave robótica em órbita, bem como testar algoritmos de controle associados à operação de manipuladores robóticos em ambiente espacial, utilizando MATLAB e Simulink como plataforma principal de desenvolvimento. A partir desse objetivo amplo, o trabalho buscou articular, em um único cenário de missão, a interação entre dinâmica orbital, movimento relativo e atuação de um manipulador montado sobre uma base flutuante.

No que se refere ao primeiro objetivo específico, voltado ao estudo de ferramentas de programação em MATLAB e Simulink, a dissertação desenvolveu um conjunto estruturado de modelos numéricos: dinâmica orbital no referencial ECI, dinâmica relativa em \gls{lvlh}, e dinâmica acoplada base--manipulador descrita por equações de Lagrange. Esses modelos foram integrados em um ambiente de simulação modular, com exportação sistemática de sinais, geração de gráficos e criação de critérios numéricos de desempenho, o que caracteriza um uso aprofundado das ferramentas de modelagem, simulação e pós-processamento disponíveis nesses ambientes.

Em relação ao segundo objetivo específico, que tratava do estudo do ambiente de atividades robóticas (incluindo aspectos de colisão, mobilidade e possíveis danos), o trabalho adotou um cenário simplificado de órbitas circulares coplanares e focou na análise de segurança em termos de janelas de posição e velocidade relativas, limites de atitude, energia cinética e momento angular. Embora não se tenha modelado explicitamente impactos, contato ou falhas estruturais, as restrições impostas às velocidades relativas e aos esforços de controle foram definidas à luz de requisitos de segurança típicos para operações de proximidade, o que permite discutir, em nível de simulação, condições mínimas para reduzir riscos de colisão e manobras abruptas.

O terceiro objetivo específico, associado ao estudo de procedimentos de otimização da mobilidade e à prevenção de falhas operacionais, foi abordado de forma indireta por meio da definição e avaliação de critérios de desempenho. Foram estabelecidas janelas admissíveis de posição e velocidade no referencial \gls{lvlh}, tolerâncias para erros de atitude, limites para energia cinética e momento angular residual, além de uma tolerância para o erro de posição da ferramenta em relação ao ponto desejado na porta de acoplamento. A comparação entre cenários com condições iniciais distintas (cenário A e cenário B) evidenciou a influência de escolhas orbitais iniciais sobre o tempo de missão e o consumo de $\Delta v$, o que se relaciona à ideia de mobilidade mais eficiente e de planejamento de missão menos suscetível a condições desfavoráveis.

Do ponto de vista das contribuições, o trabalho apresentou:
(i) uma modelagem integrada que combina dinâmica orbital no referencial ECI, movimento relativo em \gls{lvlh} e dinâmica base--manipulador em base flutuante;
(ii) uma arquitetura de controle em camadas, englobando transferência orbital, controle translacional no referencial \gls{lvlh}, controle de atitude e controle das juntas do manipulador;
(iii) critérios numéricos explícitos para avaliar posição, velocidade, atitude, energia, momento angular e erro da ferramenta; e
(iv) cenários de simulação que permitem analisar o impacto de diferentes condições iniciais sobre o tempo de missão e o custo propulsivo. Em conjunto, esses elementos atendem ao objetivo geral de estudar, em ambiente de simulação, a operação de manipuladores robóticos em órbita e fornecem uma base coerente para investigações futuras mais detalhadas sobre aproximação e ancoragem com perseguidor em base flutuante.
